% ============================================================
% main.tex — selve formelarket
%
% Denne fila er innhold. Alle typografiske valg bor i
% formelarkstil.sty, som også laster babel og siunitx.
%
% Arket er stående A4, brettet på tvers: to liggende A5-sider
% oppå hverandre. Øvre halvdel konstanter, nedre halvdel
% formler, hver halvdel i to spalter. Eksempelinnholdet er et
% formelark til Fysikk 1.
%
% Kompiler med pdflatex:
%   latexmk -pdf main
% ============================================================


\documentclass[10pt]{article}


% ============================================================
% LAST INN FORMELARKSTIL
%
% [tett]: tettere formelrader og mindre luft over
% overskriftene, for et ark med flere formler enn 10 pt gir
% plass til. Sett i tillegg \begin{halvdel}[\small]{...}{...}
% på halvdelen som er full.
% ============================================================
\usepackage{formelarkstil}
% \usepackage[tett]{formelarkstil}


% ============================================================
% EGENDEFINERTE PAKKER
%
% Legg til prosjektspesifikke pakker her.
% Ikke sett pakker andre steder i filen.
% ============================================================

% --- Vektorer: \vv{F}, flat pil over symbolet. Indekser utenfor: \vv{F}_1 ---
\usepackage[e]{esvect}


\begin{document}

% ============================================================
% ØVRE HALVDEL — KONSTANTER
% ============================================================
\begin{halvdel}{Fysikk 1}{Konstanter}
\gruppe{Konstanter}
\begin{konstanter}
Konstanten i Wiens forskyvningslov & $a$ & \qty{2.90e-3}{K.m} \\
Boltzmanns konstant       & $k$      & \qty{1.38e-23}{J/K} \\
Stefan--Boltzmann-konstanten & $\sigma$ & \qty{5.67e-8}{W/m^2K^4} \\
Plancks konstant          & $h$      & \qty{6.63e-34}{J.s} \\
Bohrs konstant            & $B$      & \qty{2.18e-18}{J} \\
Lysfarten i vakuum        & $c$      & \qty{3.00e8}{m/s} \\
Elementærladningen        & $e$      & \qty{1.60e-19}{C} \\
Tyngdeakselerasjonen      & $g$      & \qty{9.81}{m/s^2} \\
Atommasseenheten          & $u$      & \qty{1.66e-27}{kg} \\
Lysår                     & $1\,\text{ly}$ & \qty{9.46e15}{m} \\
\end{konstanter}

\gruppe{Masser}
\begin{tabellto}
Elektron     & $\num{5.4858e-4}$ u $= \qty{9.1094e-31}{kg}$ \\
Proton       & $\num{1.007276}$ u $= \qty{1.6726e-27}{kg}$ \\
Nøytron      & $\num{1.008665}$ u $= \qty{1.6749e-27}{kg}$ \\
Hydrogen     & $\num{1.007825}$ u \\
Deuterium    & $\num{2.014102}$ u \\
Tritium      & $\num{3.016049}$ u \\
Helium-3     & $\num{3.016029}$ u \\
Helium-4     & $\num{4.002603}$ u \\
Alfapartikkel & $\num{4.001506}$ u \\
\end{tabellto}

\gruppe{Jorda}
\begin{tabellto}
Masse                          & \qty{5.972e24}{kg} \\
Middelradius                   & \qty{6371}{km} \\
Middelavstand fra sola         & \qty{1.496e11}{m} \\
Lysintensitet utenfor atmosfæren & \qty{1361}{W/m^2} \\
\end{tabellto}

\gruppe{Sola}
\begin{tabellto}
Masse               & \qty{1.99e30}{kg} \\
Radius              & \qty{6.96e8}{m} \\
Overflatetemperatur & \qty{5780}{K} \\
Utstrålt effekt     & \qty{3.83e26}{W} \\
\end{tabellto}

\gruppe{SI-prefikser}
\begin{tabularx}{\columnwidth}{@{}l l Y l l Y l l l@{}}
$10^{-18}$ & atto  & a          & $10^{-3}$ & milli & m  & $10^{3}$  & kilo & k \\
$10^{-15}$ & femto & f          & $10^{-2}$ & centi & c  & $10^{6}$  & mega & M \\
$10^{-12}$ & piko  & p          & $10^{-1}$ & desi  & d  & $10^{9}$  & giga & G \\
$10^{-9}$  & nano  & n          & $10^{1}$  & deka  & da & $10^{12}$ & tera & T \\
$10^{-6}$  & mikro & \textmu{}  & $10^{2}$  & hekto & h  & $10^{15}$ & peta & P \\
\end{tabularx}
\end{halvdel}

% ============================================================
% NEDRE HALVDEL — FORMLER
% ============================================================
\begin{halvdel}{Fysikk 1}{Formler}
\gruppe{Bevegelse og krefter}
\begin{formler}
$v = at + v_0$ & $s = \dfrac{v + v_0}{2}\,t$ & \\
$v^2 - v_0^2 = 2as$ & $s = \frac{1}{2}at^2 + v_0 t$ & \\
\end{formler}
\formelrad{$\Sigma F = 0 \iff v$ konstant og rettlinjet}
\begin{formler}
$\Sigma\vv{F} = m\vv{a}$ & $F^{*} = F$ & $G = mg$ \\
$R = \mu N$ & $B = \rho_f V g$ & $L = kv^2$ \\
$\vv{p} = m\vv{v}$ & $\Sigma\vv{F}\cdot t = \Delta\vv{p}$ & \\
\end{formler}
\formelrad{$m_A v_A + m_B v_B = m_A v_{A0} + m_B v_{B0}$}

\gruppe{Mekanisk energi og effekt}
\begin{formler}
$W = Fs\cos\theta$ & $W_{\Sigma F} = \Delta E_k$ & $P = W/t$ \\
$E_k = \frac{1}{2}mv^2$ & $E_p = mgh$ & $E = E_k + E_p$ \\
\end{formler}
\formelrad{$\frac{1}{2}mv^2 + mgh = \frac{1}{2}mv_0^2 + mgh_0$}
\formelrad{$\eta = \dfrac{\text{nyttbar energi}}{\text{tilført energi}} = \dfrac{\text{nyttbar effekt}}{\text{tilført effekt}}$}

\gruppe{Trykk og temperatur}
\begin{formler}
$p = F/A$ & $T = t + 273$ & $E_k = \frac{3}{2}kT$ \\
$c = \dfrac{Q}{m\,\Delta T}$ & $C = \dfrac{Q}{\Delta T}$ & $\Delta U = W + Q$ \\
\end{formler}
\columnbreak
\gruppe{Trykk og temperatur (forts.)}
\begin{formlerto}
$\mathit{COP} = \dfrac{Q_{\text{avgitt}}}{E_{\text{elektrisk}}}$ &
$\mathit{COP}_{\text{Carnot}} = \dfrac{T_{\text{avgitt}}}{T_{\text{avgitt}} - T_{\text{opptatt}}}$ \\
\end{formlerto}

\gruppe{Bølger og stråling}
\begin{formler}
$f = 1/T$ & $v = \lambda f$ & $U = P/A$ \\
$\lambda_{\text{topp}}\cdot T = a$ & $U = \sigma T^4$ & $I = \dfrac{P}{4\pi r^2}$ \\
\end{formler}

\gruppe{Atomfysikk og kjernefysikk}
\begin{formler}
$E = hf = \dfrac{hc}{\lambda}$ & $E_n = -\dfrac{B}{n^2}$ & \\
\end{formler}
\begin{formlerto}
$hf = E_n - E_m,\quad n > m$ & $hf = B\left(\dfrac{1}{m^2} - \dfrac{1}{n^2}\right)$ \\
\end{formlerto}
\begin{formler}
$A = Z + N$ & $E_0 = mc^2$ & $\Delta E_0 = \Delta m\cdot c^2$ \\
\end{formler}

\gruppe{Elektrisitet}
\begin{formler}
$I = \dfrac{Q}{t}$ & $U = \dfrac{W}{Q}$ & $R = \dfrac{U}{I}$ \\
$W = UIt$ & $P = UI = RI^2 = \dfrac{U^2}{R}$ & \\
\end{formler}
\formelrad{Seriekobling: \quad $R_{\text{tot}} = R_1 + R_2 + \cdots + R_n$}
\formelrad{Parallellkobling: \quad $R_{\text{tot}} = \left(\dfrac{1}{R_1} + \dfrac{1}{R_2} + \cdots + \dfrac{1}{R_n}\right)^{-1}$}
\end{halvdel}

\end{document}
